\mediumtitle{Awards}

\begin{singletablelist}

	% \dtlist{Summer 2019}{\textbf{Awarded a mini-grant} (£1000) to hire an undergraduate student through the ``King's Undergraduate Research Fellowship (KURF)''.\\
% 	& }
%
% 	\dtlist{Summer 2018}{\textbf{Awarded a mini-grant} (£1000) to hire an undergraduate student through the ``King's Undergraduate Research Fellowship (KURF)''. \\
% 	& }
	
	\stlist{July 2018}{\textbf{Winner of one of the challenges at the \emph{BioDataHack 2018 - Genomic, Biodata and Improving Health Outcomes}}, a competition organised by the Wellcome Trust (2-3 July 2018, Wellcome Genome Campus, Hinxton, UK). The challenge, undertaken in collaboration with Jun Aruga, Oliver Giles, Ioannis Valasakis, and Chen Zhang, required the development of an innovative solution to answer the question \emph{``How can we use mobile technology to transform biological data processing?''}. The proposed solution was then assessed on scientific and methodological quality as well as creativity, communication skills, and ability to work as part of a team. \\
	& Alessia Visconti's team proposed the development of a device for continuous monitoring of the gut microbiome in patients with inflammatory bowel disease (IBD). The device integrated a metagenomic analysis pipeline developed by Alessia Visconti (\emph{YAMP})~\cite{Vis18b} with a shallow neural network to classify healthy individuals and patients with one of the two main forms of the disease, Crohn's disease and ulcerative colitis. The device then used a multivariate linear regression model to derive a quantitative score of disease progression, enabling the near-real-time monitoring of the effects of lifestyle changes or pharmacological treatment. During the two-day hackathon, the team successfully ported \emph{YAMP} to Arm’s 64-bit architecture, where the data could be processed within a few hours, showing that the analysis of microbial data can be successfully performed outside centralised data centres. Finally, the team showed that the proposed approach was able to separate patients with IBD from healthy controls with more than 90\% accuracy. The next step would be developing models offering personalised advice and treatments based on interventions that have been effective in patients with similar microbiome profiles.
	}

	% \dtlist{Summer 2017}{\textbf{Awarded a mini-grant} (£1000) to hire an undergraduate student through the ``King's Undergraduate Research Fellowship (KURF)''. \\
	% & }

    \stlist{Jan 2012}{\textbf{Awarded two years postdoctoral fellowship} by the Italian Minister of Education, University and Research (MIUR; duration: January 2012 - December 2013). 
	  }

	\stlist{Jan 2012}{\textbf{Awarded a postdoctoral training grant} (\euro{3,000}) by the Regione Piemonte. 
	 }

	% \dtlist{Summer 2011}{\textbf{Participation} to the \emph{``DREAM6 -- Promoter Activity Prediction Challenge''}. The proposed approach (developed with Ali Altıntas and Chris Workman) which combined results from two well-known machine learning approaches (regression trees and support vector machines for regression), ranked 8th out of 21 participants. \\
	% & }
	
	
	\stlist{Jun - Oct 2010}{\textbf{Winner of one of the challenges presented in the \emph{DREAM5 -- Network Inference Challenge}}, an international competition organised by the Dialogue for Reverse Engineering Assessments and Methods (DREAM) project to benchmark tools for the inference of gene regulatory networks from gene expression data. Methods were evaluated on simulated datasets and real data from \emph{E. coli} and \emph{S. cerevisiae}. Performance was assessed as a binary classification problem (regulatory edge present/absent) using reference (\emph{gold-standard}) networks and standard metrics (AUROC and AUPR). \\
	& The approach developed by Alessia Visconti together with Roberto Esposito and Francesca Cordero~\cite{Vis11b,Mar12}, which integrated established statistical methods for gene regulatory network inference into a Naive Bayes classifier, was compared with 35 other methods, ranking 3rd overall on real biological data and 1st in the reconstruction of the \emph{S. cerevisiae} gene regulatory network. %The results led to the inclusion of the method in the challenge's summary article~\cite{Mar12}.
	}


	\dtlist{Jan 2009}{\textbf{Awarded three years PhD fellowship} by the Italian Minister of Education, University and Research (MIUR; duration: January 2009 - December 2011). 
	}

\end{singletablelist}



